\documentclass[10pt,a4paper]{amsart}
\usepackage[margin=19mm]{geometry}
\usepackage{amsmath,amssymb,mathtools}
\usepackage[scr=rsfs,cal=euler]{mathalfa}
\usepackage{xfrac}
\usepackage[unicode,hidelinks,bookmarks=false]{hyperref}
\newcommand{\ie}{i.e.\ }
\newcommand{\cf}{cf.\ }
\newcommand{\resp}{respectively\ }
\newcommand{\Subsection}{\subsection}
\providecommand{\nobreakdash}{\nobreak\hbox{-}}
\setlength{\emergencystretch}{2em}
\multlinegap0pt
\usepackage{xcolor}
\usepackage[shortlabels]{enumitem}
\usepackage{algorithm}\usepackage{algpseudocode}
\usepackage{siunitx}
\usepackage{amssymb}
\usepackage{bm}
\usepackage{mathtools}
\newtheorem{lemma}{Lemma}
\newcommand{\E}{\mathbb{E}}
\newcommand{\R}{\mathbb{R}}
\newcommand{\abs}[1]{\left\lvert #1 \right\rvert}
\newcommand{\dd}{\,\mathrm{d}}
\newcommand{\defeq}{\triangleq}
\providecommand{\paragraph}{}
\title{Optimal Scalar Quantization for Matrix Multiplication: Closed-Form Density and Phase Transition}
\author{Calvin Ang, Sungyoon Kim, Mert Pilanci}
\date{Source: arXiv:2603.19559v1 (CC BY 4.0)}
\begin{document}
\maketitle

\noindent\emph{Source and licence.} Optimal Scalar Quantization for Matrix Multiplication: Closed-Form Density and Phase Transition. Version arXiv:2603.19559v1 (CC BY 4.0), distributed by arXiv under the Creative Commons Attribution 4.0 International licence, http://creativecommons.org/licenses/by/4.0/, as recorded in the arXiv licence metadata for this exact version and shown on the abstract page https://arxiv.org/abs/2603.19559v1. Full licence notice: \texttt{licenses/arxiv-2603.19559-CC-BY-4.0.txt}.\par\smallskip

\noindent\textit{Editorial scope.} Counted unit: the article's lemma lem:mixed\_cell\_bound (source0.tex lines 915--928). The article's complete original proof of it is delivered with it (source0.tex lines 929--1014, one \texttt{proof} environment of 86 lines). No mathematics is written, repaired or re-derived: the statement and the proof are the article's own lines, in the article's own notation, and no retained line is substituted or re-flowed.  Retained uncounted as the source-local context the proof needs: lines 726--741; lines 732--739; lines 791--799.  Nothing was omitted from any retained span, so the package carries no omission record.  No retained line contains a citation, so the wrapper carries no bibliography and no bibliography entry.  No retained line contains a figure environment, an image-inclusion command or drawing code, so no image is delivered and the package declares no asset dependency.  Editorial formatting only, in addition: an AMS \texttt{amsart} preamble that restates the article's own declarations and macros used by the retained lines, namely the environments \texttt{lemma} on separate counters, together with the standard packages those lines need.  The retained environments print in this document as Lemma 1 in order of appearance, whereas the article prints the same items with the numbering of its own full document.  The complete native source of the article as distributed in the arXiv e-print for this exact version is bundled unchanged as \texttt{sources/196785/source0.tex}.
\begin{lemma}[Reproduction points are second-order close to cell midpoints]\label{lem:r_midpoint}
Assume $\lambda$ is continuously differentiable. Then for each fixed compact interval $[-M,M]$ there exists $C_M<\infty$ and $K_M$ such that for all $K\ge K_M$ and all cells $I_i$ intersecting $[-M,M]$,
\begin{equation}\label{eq:r_minus_m}
\abs{r_i-m_i}\le C_M\,\Delta_i^2.
\end{equation}
Consequently,
\begin{equation}\label{eq:first_second_moment_id}
\begin{aligned}
\int_{x_{i-1}}^{x_i} (x-r_i)\, \mathrm{d}x
&= \Delta_i(m_i-r_i)=O(\Delta_i^3), \\
\int_{x_{i-1}}^{x_i} (x-r_i)^2\, \mathrm{d}x
&= \frac{\Delta_i^3}{12}+O(\Delta_i^5).
\end{aligned}
\end{equation}
uniformly over such cells.
\end{lemma}

\begin{equation}
\begin{aligned}
\int_{x_{i-1}}^{x_i} (x-r_i)\, \mathrm{d}x
&= \Delta_i(m_i-r_i)=O(\Delta_i^3), \\
\int_{x_{i-1}}^{x_i} (x-r_i)^2\, \mathrm{d}x
&= \frac{\Delta_i^3}{12}+O(\Delta_i^5).
\end{aligned}
\end{equation}

\begin{lemma}[A weighted first-moment bound for companders]\label{lem:first_moment}
Let $X$ have density $f$ and let $g:\R\to\R$ be continuously differentiable. Let $Q_{K,\lambda}$ be a companding quantizer with continuously differentiable point density $\lambda>0$.
Assume that $q(x)\defeq g(x)f(x)$ is twice continuously differentiable.
Then
\begin{equation}\label{eq:first_moment_order}
\E\big[g(X)\,(X-Q_{K,\lambda}(X))\big]=O\!\left(\frac{1}{K^2}\right),
\end{equation}
as $K\to\infty$.
\end{lemma}

\begin{lemma}[A generic mixed-cell bound]\label{lem:mixed_cell_bound}
Let $\phi:\R^2\to\R$ be twice continuously differentiable. Consider companding quantizers $Q_X=Q_{K_X,\lambda_X}$ and $Q_Y=Q_{K_Y,\lambda_Y}$ with continuosly differentiable point densities. Let $r_i$ and $s_j$ denote the reproduction points and let $\Delta_i$ and $\delta_j$ denote the corresponding cell widths for $Q_X$ and $Q_Y$ respectively.
Assume the function
\[
M_\phi(x,y)\defeq \max_{|\alpha|+|\beta|=2}\abs{\partial_x^\alpha\partial_y^\beta \phi(x,y)}
\]
is integrable with respect to $f_{X,Y}(x,y)\dd x\dd y$ when weighted by $(\lambda_X(x)\lambda_Y(y))^{-2}$ on compact sets (as made explicit in the proof below).
Then, as $K_X,K_Y\to\infty$,
\begin{equation}\label{eq:mixed_term_order_template}
\sum_{i=1}^{K_X}\sum_{j=1}^{K_Y}\int_{I_i}\int_{J_j}
\phi(x,y)\,(x-r_i)(y-s_j)\dd y\dd x
=O\!\left(\frac{1}{K_X^2K_Y^2}\right).
\end{equation}
\end{lemma}

\begin{proof}
Let $I_i=[x_{i-1},x_i)$ and $J_j=[y_{j-1},y_j)$ be the cells in $x$ and $y$. Fix $M>0$ and split the double sum into (i) rectangles intersecting $[-M,M]^2$ and (ii) rectangles in the complement. We bound both contributions.

\paragraph*{Step 1: Rectangles intersecting a compact set.}
For a rectangle $I_i\times J_j$ intersecting $[-M,M]^2$, pick the center point $(r_i,s_j)$ and perform a second-order Taylor expansion of $\phi$ around $(r_i,s_j)$:

\begin{align*}
    \phi(x,y) &=\phi(r_i,s_j)+\phi_x(r_i,s_j)(x-r_i)\\
    &+\phi_y(r_i,s_j)(y-s_j)+R_{i,j}(x,y),
\end{align*}
where the remainder satisfies
\[
\abs{R_{i,j}(x,y)}
\le \frac{1}{2} M_\phi(\zeta_{i,j}(x,y))\bigl((x-r_i)^2+(y-s_j)^2\bigr)
\]
for some $\zeta_{i,j}(x,y)\in I_i\times J_j$.

Multiply the Taylor expansion by $(x-r_i)(y-s_j)$ and integrate over $I_i\times J_j$. The constant term contributes
\[
\phi(r_i,s_j)\left(\int_{I_i}(x-r_i)\dd x\right)\left(\int_{J_j}(y-s_j)\dd y\right).
\]
By Lemma~\ref{lem:r_midpoint}, each one-dimensional integral is $O(\Delta_i^3)$ and $O(\delta_j^3)$ on a compact region, hence the constant term contribution is $O(\Delta_i^3\delta_j^3)$.

The $\phi_x$ term contributes
\[
\phi_x(r_i,s_j)\left(\int_{I_i}(x-r_i)^2\dd x\right)\left(\int_{J_j}(y-s_j)\dd y\right)
=O(\Delta_i^3\delta_j^3),
\]
since $\int_{I_i}(x-r_i)^2\dd x=O(\Delta_i^3)$ by \eqref{eq:first_second_moment_id}. The $\phi_y$ term is analogous.

For the remainder term, note that on $I_i\times J_j$ we have $\abs{x-r_i}=O(\Delta_i)$ and $\abs{y-s_j}=O(\delta_j)$, so
\[
\abs{R_{i,j}(x,y)(x-r_i)(y-s_j)}
\le C\, M_\phi(\zeta_{i,j}(x,y))\bigl(\Delta_i^3\delta_j+\Delta_i\delta_j^3\bigr).
\]
Integrating over $I_i\times J_j$ contributes at most

\begin{align*}
C\left(\Delta_i^3\delta_j+\Delta_i\delta_j^3\right)&\int_{I_i}\int_{J_j} M_\phi(\zeta_{i,j}(x,y))\dd y\dd x
\le\\
&C'\Delta_i\delta_j\left(\Delta_i^2+\delta_j^2\right)\sup_{I_i\times J_j}M_\phi.
\end{align*}

On compact sets, $\sup_{I_i\times J_j}M_\phi$ is bounded and $\Delta_i,\delta_j=O(1/K_X),O(1/K_Y)$, so this remainder contribution is also $O(\Delta_i^3\delta_j^3)$.

Combining all pieces, we have shown that for rectangles intersecting $[-M,M]^2$,
\[
\int_{I_i}\int_{J_j}\phi(x,y)(x-r_i)(y-s_j)\dd y\dd x = O(\Delta_i^3\delta_j^3),
\]
uniformly. Summing over the $O(K_XK_Y)$ rectangles intersecting $[-M,M]^2$ and using $\Delta_i=O(1/K_X)$, $\delta_j=O(1/K_Y)$ on compacts yields a total compact contribution of order $O(K_XK_Y\cdot K_X^{-3}K_Y^{-3})=O(K_X^{-2}K_Y^{-2})$.

\paragraph*{Step 2: Tail rectangles.}
On the complement of $[-M,M]^2$ we bound the entire integral by absolute values. Since $\abs{x-r_i}\le \abs{x}+\abs{r_i}$ and similarly for $y-s_j$, and since $(X,Y)$ has finite $(4+\epsilon)$ moments, by choosing $M$ large we can make the tail probability $\Pr((X,Y)\notin[-M,M]^2)$ arbitrarily small. The detailed bound follows the same truncation logic as in Lemma~\ref{lem:first_moment}: first control the tail by moments and then let $M\to\infty$. This yields a tail contribution that is $o(K_X^{-2}K_Y^{-2})$.

\paragraph*{Step 3: Combine and conclude.}

Let the rectangles which intersect $[-M,M]^2$ be $C := \{(i,j) : I_i \times J_j \cap [-M,M]^2 \neq \emptyset  \}$ and $\psi_{i,j}(x,y) := \phi(x,y)\,(x-r_i)(y-s_j)$. Then

\begin{align*}
    \sum_{i=1}^{K_X}\sum_{j=1}^{K_Y}\int_{I_i}\int_{J_j}\psi_{i,j}(x,y)
\dd y\dd x &= \sum_{(i,j) \in C} \int_{I_i}\int_{J_j}  
\psi_{i,j}(x,y)\dd y\dd x\\
&+  \sum_{(i,j) \not\in C} \int_{I_i}\int_{J_j} \psi_{i,j}(x,y) \dd y \dd x\\
\end{align*}

and from Step 1 we know that $\sum_{(i,j) \in C} \int_{I_i}\int_{J_j}  
\psi_{i,j}(x,y)\dd y\dd x = O(K_X^{-2} K_Y^{-2})$. 

From Step 2 we know that

\begin{align*}
    \int_{I_i}\int_{J_j} \psi_{i,j}(x,y) \dd y \dd x &= o(K_X^{-2} K_Y^{-2}), ~~ \forall (i,j) \not\in C\\
    \implies \sum_{(i,j) \not\in C} \int_{I_i}\int_{J_j} \psi_{i,j}(x,y) \dd y \dd x &= O(K_X K_Y) o(K_X^{-2} K_Y^{-2})\\
    &= O(K_X^{-2} K_Y^{-2}).
\end{align*}


Thus we see that

\begin{align*}
    \sum_{i=1}^{K_X}\sum_{j=1}^{K_Y}\int_{I_i}\int_{J_j}\psi_{i,j}(x,y)
\dd y\dd x &= O(K_X^{-2} K_Y^{-2})
\end{align*}

which is exactly \eqref{eq:mixed_term_order_template}.
\end{proof}

\end{document}
